\setcounter{ex}{0}
\section{CẤP SỐ CỘNG - CẤP SỐ NHÂN}
\subsection{Bài tập rèn luyện}
\Opensolutionfile{ans}[ans/DA-CD6-CAP-SO-BTRL]
\begin{ex}%[1D2H3-3]
 Cho cấp số nhân $(u_n)$ với $u_1=1$; $q=\dfrac{1}{10}$. Hỏi $\dfrac{1}{10^{103}}$ là số hạng thứ bao nhiêu?
 \choice
 {$105$}
 {$103$}
 {\True $104$}
 {$102$}
 \loigiai{
 Ta có số hạng thứ $n$ là $u_n=u_1 \cdot q^{n-1}=1 \cdot \left(\dfrac{1}{10}\right)^{n-1}=\dfrac{1}{10^{n-1}}$.\\
 Mà $u_n=\dfrac{1}{10^{103}}$ suy ra $n-1=103$ hay $n=104$.
 }
\end{ex}
\begin{ex}%[1D2N2-2]
 Một cấp số cộng gồm $5$ số hạng. Hiệu số hạng đầu và số hạng cuối bằng $20$. Tìm công sai $d$ của cấp số cộng đã cho.
 \choice
 {\True $d=-5$}
 {$d=4$}
 {$d=-4$}
 {$d=5$}
 \loigiai
 {
 Gọi năm số hạng của cấp số cộng đã cho là $u_1; u_2; u_3; u_4; u_5$.\\
 Theo đề bài ta có $u_1 - u_5 = 20 \Leftrightarrow u_1 - \left(u_1+4d\right)=20\Leftrightarrow d=-5$.
 }
\end{ex}
\begin{ex}%[1D2N2-2]
 Cho cấp số cộng $\left(u_n \right)$ có $u_2=4$; $u_4=2$. Tính số hạng đầu tiên $u_1$ và công sai $d$ của cấp số cộng.
 \choice
 {$u_1=1$ và $d=1$}
 {$u_1=6$ và $d=1$}
 {\True $u_1=5$ và $d=-1$}
 {$u_1=-1$ và $d=-1$}
 \loigiai{
 Ta có $\heva{&{u_2=4} \\
 &{u_4=2}} \Leftrightarrow \heva{&{u_1+d=4} \\
 &{u_1+3d=2}} \Leftrightarrow \heva{&{u_1=5} \\
 &{d=-1}.}$ \\ Vậy cấp số cộng đã cho có số hạng đầu tiên $u_1=5$ và công sai $d=-1$.
 }
\end{ex}
\begin{ex}%[1D2N2-2]
 Trong các dãy số $(u_n)$ với số hạng tổng quát sau, dãy số nào là cấp số cộng?
 \choice
 {$u_n=5^n$}
 {\True $u_n=2-5n$}
 {$u_n=5^n-2$}
 {$u_n=5+n^2$}
% \loigiai{
% Xét $u_n=2-5n \Rightarrow u_{n+1}=2-5(n+1)$.\\
% Ta có $u_{n+1}-u_n=2-5(n+1)-(2-5n)=-5$ (hằng số).\\
% Vậy $u_n=2-5n$ là cấp số cộng có số hạng đầu $u_1=2$ và công sai $d=-5$.}
\end{ex}
\begin{ex}%[1D2N2-2]
 Cho cấp số cộng $(u_n)$ có $u_1=2$, $u_9=34$. Tìm công sai $d$ của cấp số cộng $(u_n)$ ?
 \choice
 {\True $d=4$}
 {$d=-4$}
 {$d=3$}
 {$d=-3$}
 \loigiai{
 Ta có $u_9=u_1+8d
 \Leftrightarrow 34=2+8d
 \Leftrightarrow d=4$.\\
 Vậy cấp số cộng có công sai $d=4$.
 }
\end{ex}
\begin{ex}%[1D2N2-3]
 Cho cấp số cộng $(u_n)$ biết $u_1=-1$; $d=2$; $u_n=43$. Hỏi cấp số cộng đó có bao nhiêu số hạng?
 \choice
 {$20$}
 {\True $23$}
 {$22$}
 {$21$}
 \loigiai{ Ta có
 $u_n=u_1+(n-1)d \Leftrightarrow 43=-1+(n-1)\cdot2\Leftrightarrow n=23$.
 }
\end{ex}
\begin{ex}%[1D2N2-3]
 Cho cấp số cộng $(u_n)$ với $u_1=4$ và công sai $d=-3$. Giá trị của $u_5$ bằng
 \choice
 {$16$}
 {$19$}
 {\True $-8$}
 {$-11$}
 \loigiai{
 Cấp số cộng $(u_n)$ có $u_5=u_1+4d=4+4\cdot (-3)=-8$.
 }
\end{ex}
\begin{ex}%[1D2N2-4]
 Cho cấp số cộng $(u_n)$ có số hạng đầu $u_1=2$ và công sai $d=3$. Giá trị $u_{2024}$ bằng
 \choice
 {$6065$}
 {$6068$}
 {\True $6071$}
 {$6074$}
 \loigiai{
 Áp dụng công thức số hạng tổng quát
 $u_{2024} =u_1 + 2023d = 2 + 2023 \cdot 3 =6071$.
 }
\end{ex}
\begin{ex}%[1D2H2-2]
 Cho một cấp số cộng $\left(u_n\right)$ có $u_1=5$ và tổng của 40 số hạng đầu là $3320$. Tìm công sai của cấp số cộng đó.
 \choice
 {$-4$}
 {$8$}
 {$-8$}
 {\True $4$}
 \loigiai{
 Gọi $d$ là công sai của cấp số cộng.\\
 Ta có tổng 40 số hạng đầu của cấp số cộng là
 \begin{eqnarray*}
 & & S_{40}=\dfrac{40\left(2u_1+39d\right)}{2}=3320\\
 &\Leftrightarrow & \dfrac{40\left(2.5+39d\right)}{2}=3320\\
 &\Leftrightarrow & d=4.
 \end{eqnarray*}
}\end{ex}
\begin{ex}%[1D2H2-2] :
 Cho một cấp số cộng có số hạng tổng quát là $ u_{n}=3+4 n$. Công sai của cấp số cộng này bằng
 \choice
 {$3$}
 {$7$}
 {$11$}
 {\True $4$}
 \loigiai{
 Ta có $u_1= 3+4\cdot 1=7; u_2= 3+4\cdot 2=11$. Suy ra $d=u_2-u_1=11-7=4$.
 }
\end{ex}
\begin{ex}%[1D2H2-2]
 Trong các dãy số nào sau đây, dãy số nào là cấp số cộng?
 \choice
 {$u_n=\dfrac{1}{n}$}
 {\True $u_n=u_{n-1}-2$, $\forall n \ge 2$}
 {$u_n=2^n-1$}
 {$u_n=2u_{n-1}$, $\forall n \ge 2$}
 \loigiai{Xét dãy số $(u_n)$ với $u_n=u_{n-1}-2$, $\forall n \ge 2$.\\
 Ta có $u_n-u_{n-1}=-2$, $\forall n \ge 2$ nên $(u_n)$ là một cấp số cộng.
 }
\end{ex}
\begin{ex}%[1D2H2-2]
 Xác định số hạng đầu $u_1$ và công sai $d$ của cấp số cộng $(u_n)$ có $u_9=5u_2$ và \\$u_{13}=2u_6+5$.
 \choice
 {\True $u_1=3$ và $d=4$}
 {$u_1=3$ và $d=5$}
 {$u_1=4$ và $d=5$}
 {$u_1=4$ và $d=3$}
 \loigiai{Ta có $\heva{&u_9=5u_2\\&u_{13}=2u_6+5} \Rightarrow \heva{&u_1+8d=5(u_1+d)\\&u_1+12d=2(u_1+5d)+5} \Rightarrow \heva{&4u_1-3d=0\\&u_1-2d=-5} \Rightarrow \heva{&u_1=3\\&d=4.}$\\
 Vậy $u_1=3$ và $d=4$.
 }
\end{ex}
\begin{ex}%[1D2H2-2]
 Cho một cấp số cộng $(u_n)$ có $u_1=5$ và tổng của $40$ số hạng đầu là $3\,320$. Tìm công sai của cấp số cộng $(u_n)$.
 \choice
 {$-4$}
 {$8$}
 {$-8$}
 {\True $4$} \loigiai{Gọi $d$ là công sai của cấp số cộng.\\ Ta có tổng $40$ số hạng đầu của cấp số cộng là
 \allowdisplaybreaks
 \begin{align*}
 S_{40}=\dfrac{40(2u_1+39d)}{2}=3\,320
 \Leftrightarrow \dfrac{40(2\cdot 5+39d)}{2}=3\,320 \Leftrightarrow d=4.
 \end{align*}
 }
\end{ex}
\begin{ex}%[1D2H2-2]
 Cho cấp số cộng $(u_n)$ biết $u_5=18$ và $4S_n=S_{2n}$. Số hạng đầu $u_1$ và công sai $d$ của cấp số cộng là
 \choice
 {$u_1=3$, $d=2$}
 {$u_1=2$, $d=3$}
 {$u_1=2$, $d=2$}
 {\True $u_1=2$, $d=4$}
 \loigiai{
 Từ giả thiết ta có hệ
 \begin{eqnarray*}
 &&\heva {&u_5=18\\&4S_n=S_{2n}}\\
 &\Leftrightarrow & \heva {&u_1+4d=18\\&4n\cdot \dfrac{u_1+u_n}{2}=2n\cdot \dfrac{u_1+u_{2n}}{2}}\\&\Leftrightarrow& \heva {&u_1+4d=18\\&2(u_1+u_n)=u_1+u_{2n}}\\
 &\Leftrightarrow & \heva {&u_1+4d=18\\&2\left[2u_1+(n-1)d\right]=2u_1+(2n-1)d}\\&\Leftrightarrow& \heva {&u_1+4d=18\\&2u_1-d=0}\\&\Leftrightarrow& \heva {&u_1=2\\&d=4.}
 \end{eqnarray*}
 }
\end{ex}
\begin{ex}%[1D2H2-2]
 Cho cấp số cộng $\left(u_n\right)$ với $u_{17}=33$ và $u_{33}=65$ thì công sai bằng
 \choice
 {$1$}
 {$3$}
 {$-2$}
 {\True $2$}
 \loigiai{
 Gọi $u_1$, $d$ lần lượt là số hạng đầu và công sai của cấp số cộng $\left(u_n\right)$.\\
 Khi đó, ta có $u_{17}=u_1+16d$, $u_{33}=u_1+32d$.\\
 Suy ra $u_{33}-u_{17}=65-33\Leftrightarrow 16d=32\Leftrightarrow d=2$.\\
 Vậy công sai bằng $2$.}
\end{ex}
\begin{ex}%[1D2H2-3]
 Cho cấp số cộng $(u_n)$ với $u_1=2$; $d=9$. Khi đó số $2018$ là số hạng thứ mấy trong dãy?
 \choice
 {$226$}
 {\True $225$}
 {$223$}
 {$224$}
 \loigiai{ Ta có
 $u_n=u_1+(n-1)d \Leftrightarrow 2018=2+(n-1 )\cdot9\Leftrightarrow n=225$.
 }
\end{ex}
\begin{ex}%[1D2H2-6]
 Cho dãy $(a_n)$ là một cấp số cộng, biết $a_3+a_8+a_{10}+a_{16}+a_{18}+a_{23}=126$. Tính tổng của $25$ số hạng đầu tiên của dãy $(a_n)$.
 \choice{$315 $}
 {$ 550$}
 {$ 552$}
 {\True $525 $}
 \loigiai{ Gọi $d$ là công sai của cấp số cộng.\\
 Ta có $(a_3+a_{23})+(a_8+a_{18})+(a_{10}+a_{16})=6a_{13}\Rightarrow a_{13}=21.$ \\
 Mặt khác $S_{25}=\dfrac{a_1+a_{25}}{2}\cdot 25=\dfrac{2a_1+24d}{2}\cdot 25= 25\cdot (a_1+12d)=25\cdot a_{13}=525$.
 }
 \dotlineEX{1}% Tạo các dòng kẻ trong tài liệu
\end{ex}
\begin{ex}%[1D2H2-6]
 Viết thêm năm số vào giữa số $2$ và số $20$ để được một cấp số cộng. Tổng của năm số đó là
 \choice{$77$}
 {$40$}
 {\True $55$}
 {$60$}
 \loigiai{
 Ta có $u_1=2$ và $u_7=20$. Suy ra $u_1+6d=20\Rightarrow d=3$.\\
 Khi đó $5$ số được viết thêm là $5, 8, 11, 14, 17$.\\
 Tổng của $5$ số đó là $5+8+11+14+17=55$
 }
\end{ex}
\begin{ex}%[1D2N3-2]
 Cho cấp số nhân $(u_n)$ có $u_1=-\dfrac{1}{2}$, $u_7=-32$. Tìm công bội $q$.
 \choice
 {$q=\pm\dfrac{1}{2}$}
 {\True $q=\pm 2$}
 {$q=\pm 4$}
 {$q=\pm 1$}
 \loigiai
 {
 Ta có $u_7=u_1q^6\Rightarrow -32=-\dfrac{1}{2}\cdot q^6\Rightarrow q=\pm 2$.
 }
\end{ex}
\begin{ex}%[1D2N3-2]
 Dãy số nào trong các dãy số $(u_n)$ được cho sau đây là cấp số nhân?
 \choice
 {\True $\heva{&u_1= 3 \\ &u_{n+1}= -\dfrac{u_n}{5}}$}
 {$\heva{&u_1= 1, u_2= \sqrt{2} \\ &u_{n+2}= u_{n+1} \cdot u_n}$}
 {$\heva{&u_1= 3 \\ &u_{n+1}= n \cdot u_n}$}
 {$u_n= 2 \cdot n^2$}
 \loigiai{
 Từ công thức định nghĩa cấp số nhân $u_n^2= u_{n-1} \cdot u_{n+1}$ và $u_n= u_1 \cdot q^{n-1}$ có công bội $q$ là hằng số không đổi.\\
 Đáp án \lq\lq $\heva{&u_1= 3 \\ &u_{n+1}= -\dfrac{u_n}{5}}$\rq\rq\, có $\dfrac{u_{n+1}}{u_n}= \dfrac{-1}{5}$ đúng định nghĩa nên nhận.\\
 Đáp án \lq\lq$\heva{&u_1= 1, u_2= \sqrt{2} \\ &u_{n+2}= u_{n+1} \cdot u_n}$\rq\rq, \lq\lq $\heva{&u_1= 3 \\ &u_{n+1}= n \cdot u_n}$\rq\rq\, công thức đã sai với công thức định nghĩa nên loại.\\
 Đáp án \lq\lq $u_n= 2 \cdot n^2$\rq\rq\, có $\dfrac{u_{n+1}}{u_n}= \dfrac{2 \cdot (n+1)^2}{2 \cdot n^2}= 1+ \dfrac{2n+1}{n^2}$ không là hằng số nên cũng loại.
 }
\end{ex}
\begin{ex}%[1D2N3-2]
 Một cấp số nhân có sáu số hạng, số hạng đầu bằng $2$ và số hạng thứ sáu bằng $486$. Tìm công bội $q$ của cấp số nhân đã cho.
 \choice
 {\True $ q=3$}
 {$ q=-3$}
 {$ q=2$}
 {$ q=-2$}
 \loigiai{
 Giả sử $(u_n)$ là cấp số nhân đã cho.\\
 Ta có $u_1=2$, $u_6=486$ và số hạng tổng quát $u_n=u_1\cdot q^{n-1}$. Do đó
 \begin{align*}
 u_6=u_1\cdot q^{6-1} \Leftrightarrow 486=2\cdot q^{6-1} \Leftrightarrow q^{5}=243 \Leftrightarrow q=3.
 \end{align*}
 }
 \dotlineEX{1}% Tạo các dòng kẻ trong tài liệu
\end{ex}
\begin{ex}%[1D2N3-2]
 Nếu cấp số nhân $\left(u_n\right)$ có các số hạng $u_3=2$, $u_4=-4$ thì công bội $q$ bằng bao nhiêu?
 \choice
 {$q= -1$}
 {\True $q= -2$}
 {$q=\pm 2$}
 {$q= 2$}
 \loigiai{
 Vì $\left(u_n\right)$ là cấp số nhân nên $u_4=u_3 \cdot q \Leftrightarrow q=\dfrac{u_4}{u_3}=\dfrac{-4}{2}=-2$.
 }
\end{ex}
\begin{ex}%[1D2N3-2]
 Cho cấp số nhân $(u_n)$, biết $u_1=4$ và $u_4=64$. Tính công bội $q$ của cấp số nhân đã cho.
 \choice
 {\True $q=4$}
 {$q=-4$}
 {$q=21$}
 {$q=2\sqrt{2}$}
 \loigiai
 {Ta có $u_4=64$ suy ra $u_1\cdot q^3=64$. Vậy $q=4$.
 }
\end{ex}
\begin{ex}%[1D2N3-3]%[Anh Duy]
 Cho cấp số nhân có $u_1=4$ và $q=2$. Công thức tổng quát của cấp số nhân đã cho là
 \choice
 {$u_n=4\cdot 2^n$}
 {\True $u_n=2^{n+1}$}
 {$u_n=2^{n-1}$}
 {$u_n=2\cdot 2^{n-1}$}
 \loigiai{
 Công thức tổng quát của cấp số nhân đã cho là
 \[u_n=u_1\cdot q^{n-1}=4\cdot 2^{n-1}=2^2\cdot 2^{n-1}=2^{n+1}.\]
 }
\end{ex}
\begin{ex}%[1D2H3-2]%[Dự án 11 HVA 2024-2025]%[Ngô Tất Thành]
 Tìm công bội $q$ của một cấp số nhân $(u_n)$ có $u_1=\dfrac{1}{2}$ và $u_6=16$.
 \choice
 {$q=\dfrac{1}{2}$}
 {$q=-2$}
 {\True $q=2$}
 {$q=-\dfrac{1}{2}$}
 \loigiai{
 Ta có $u_6=u_1\cdot{q^5} \Rightarrow{q^5}=\dfrac{u_6}{u_1}=\dfrac{16}{\dfrac{1}{2}}=32 \Rightarrow q=2$.}
\end{ex}
\begin{ex}%[1D2H3-4]%[Dự án 11 HVA 2024-2025]%[Vĩnh Tín]
 Cho cấp số nhân $(u_n)$ có $u_1=3$ và $q=-2$. Số $-96$ là số hạng thứ mấy của cấp số nhân đã cho?
 \choice
 {Số hạng thứ $5$}
 {\True Số hạng thứ $6$}
 {Số hạng thứ $7$}
 {Không là số hạng của cấp số đã cho}
 \loigiai{
 Ta có $-96=u_n=u_1q^{n-1}=3\cdot (-2)^{n-1}\Leftrightarrow(-1)^{n-1} \cdot 2^{n-1}=-32=(-1)^5\cdot 2^5\Leftrightarrow n=6$.}
\end{ex}
\begin{ex}%[1D2H3-5]
 Xác định $x$ để $3$ số $x-2; x+1; 3-x$ theo thứ tự lập thành một cấp số nhân:
 \choice
 {\True Không có giá trị nào của $x$}
 {$x=\pm 1$}
 {$x=2$}
 {$x=-3$}
 \loigiai{
 Ba số $x-2; x+1; 3-x$ theo thứ tự lập thành một cấp số nhân khi
 \[(x-2)(3-x)=(x+1)^2 \Leftrightarrow 2x^2-3x+7=0.\]
 }
\end{ex}
\begin{ex}%[1-TK-HK1-CT-2-2425]%[VN-MT-9-CT-10-11, Nguyễn Khánh Trọng]%[1D2H3-5]
 Tìm giá trị của $a$ để $-\dfrac{1}{5}$; $a$; $-\dfrac{1}{125}$ theo thứ tự lập thành cấp số nhân.
 \choice
 {\True $a=\pm \dfrac{1}{25}$}
 {$a=\dfrac{1}{25}$}
 {$a=\pm \dfrac{1}{5}$}
 {$a=\dfrac{1}{5}$}
 \loigiai{
 Ta có $a^2=\left( -\dfrac{1}{5} \right)\cdot\left( -\\dfrac{1}{125} \right)=\dfrac{1}{625}\Leftrightarrow a=\pm \dfrac{1}{25}$.
 }
\end{ex}
\begin{ex}%[1D2H3-5]
 Cho tam giác $ABC$ có ba góc nhọn $A$, $B$, $C$ theo thứ tự lập thành một cấp số nhân với công bội $q=2$. Tính số đo góc $A$.
 \choice
 {$\dfrac{\pi}{2}$}
 {\True $\dfrac{\pi}{7}$}
 {$\dfrac{2\pi}{7}$}
 {$\dfrac{4\pi}{7}$}
 \loigiai{
 Theo giả thiết ta có $A$, $B$, $C$ lập thành cấp số nhân với $q=2$ nên $B=2A$; $C=4A$.\\
 Mà $ABC$ là một tam giác nên $A+B+C=\pi \Leftrightarrow A+2A+4A=\pi \Leftrightarrow A=\dfrac{\pi}{7}$.
 }
\end{ex}
\begin{ex}%[1D2H3-6]
 Cho cấp số nhân $(u_n)$ có hạng đầu $u_1=2$ và tổng của $8$ số hạng đầu tiên $S_8=6560$. Tìm công bội $q$ của cấp số nhân đã cho.
 \choice
 {\True $q=3$}
 {$q=-3$}
 {$q=\dfrac{1}{3}$}
 {$q=\pm 3$}
% \loigiai{
% Ta có $S_n=\dfrac{u_1(q^n-1)}{q-1}$ suy ra $S_8=6560=\dfrac{2(q^8-1)}{q-1}\Leftrightarrow 3280=q^7+q^6+q^5+q^4+q^3+q^2+q+1$
% \[\Leftrightarrow (q-3)\left[(q^3+2q^2)^2+(3q^2+\dfrac{20}{3}q)^2+\dfrac{689}{9}q^2+364q+1093\right]=0,\]
% từ đó suy ra $q=3$.
% }
\end{ex}
\Closesolutionfile{ans}